%% toc: P15 164 164.10
\chapter{Uniqueness, Cox's Postulates and Non-Additive Belief}
\label{sec:p15-w02b-p15-032}

Subjective probability is often defended as the only rational form for
graded belief. Dubois and Prade's 2021 survey questions that claim from
several directions, and \cref{sec:p15-rationality-4-7} gives its overall
organization. This chapter supplies the arguments that the survey states
briefly and leaves unproved: the route from fair prices to one distribution,
the hidden content of Cox's postulates, the difference between lower
probabilities and belief functions, the two decision rules for imprecise
previsions and their behaviour on Ellsberg's urn, the semantic content of
the epistemic axioms, the representation of convex sets of probabilities,
and the exact class of probabilities that preserve adjunction
\cite[pp.~305--316]{DuboisPrade2021RepresentingBelief}.

The inputs are a finite nonempty state space \(S\), events \(A\subseteq S\),
confidence functions on events, sets of probability vectors on \(S\), and
real-valued gambles on \(S\). Each formal result is proved in
\cref{sec:annex-p15-w02b-p15-032}. The survey's claims are of four kinds,
and each one is tagged with its kind below: historical attributions, normative
proposals, formal statements, and reports of empirical findings. A formal
result proved here shows only what follows from its hypotheses. It does
not settle whether an agent ought to adopt those hypotheses.

\section{Historical claims as the survey reports them}
\label{sec:p15-w02b-p15-032-history}

The survey places its formal material inside a history of two traditions.
\Cref{tab:p15-w02b-p15-032-history} lists the historical claims that
\cref{sec:p15-rationality-4-7} does not already present. Each row says what
the survey asserts, where it asserts it, and what the report does not
establish. In every row the survey is the only source read here. Where the
survey itself relies on a secondary history, the table says so.

\begin{longtable}{@{}p{.25\linewidth}p{.08\linewidth}p{.19\linewidth}p{.38\linewidth}@{}}
\caption{Historical claims reported by Dubois and Prade, with their type.}
\label{tab:p15-w02b-p15-032-history}\\
\toprule Claim as reported & Page & Type & What the report does not establish \\
\midrule\endfirsthead
\caption[]{Historical claims reported by Dubois and Prade (continued)}\\
\toprule Claim as reported & Page & Type & What the report does not establish \\
\midrule\endhead
\midrule\multicolumn{4}{r}{\emph{continued on next page}}\\
\endfoot
\endlastfoot
Hintikka's epistemic logic treats knowledge as, in essence, true belief,
whereas a merely believed proposition need not be true. Ignorance about
\(p\) is expressed by believing neither \(p\) nor its negation. Knowing
each of two propositions gives knowledge of their conjunction, which
Kyburg calls adjunction.
& p.~305 & historical attribution about \cite{Hintikka1962} and
\cite{R47Kyburg1997} &
Hintikka's own axioms and semantics are not read here. The survey's
minimal logic MEL, a different system, is analysed formally in
\cref{thm:p15-w02b-p15-032-mel}. \\
Bernoulli and Lambert did not take degrees of belief to be additive.
& p.~306 & historical attribution, secondhand, through
\cite{R47Shafer1978} &
No text by Bernoulli or Lambert is read. The claim rests on Shafer's
history as the survey reports it. \\
In the mid-twentieth century de Finetti obtained additivity by reading
\(P(A)\) as a fair price for a ticket paying one unit if \(A\) is true,
with sure loss to be avoided. & p.~306 & historical attribution with a
formal core & The formal core is \cref{thm:p17-finite-price-coherence}. Whether
the fair-price commitment is rationally required is a separate normative
question; see \cref{sec:p15-w02b-p15-032-uniqueness}. \\
Later work in artificial intelligence used Cox's result to justify
probability with Bayes' rule as the only reasonable way to represent and
revise numerical belief. The survey names
\cite{R47HorvitzHeckermanLanglotz1986}, \cite{R47Cheeseman1988} and
\cite{R47Jaynes2003}. & p.~310 & historical report about a literature &
What each named work argued is not read here. The phrase ``ad nauseam''
is the survey's evaluation, not a count of publications. \\
Shafer's theory of evidence may be viewed either as one interpretation of
Dempster's 1967 upper and lower probabilities for imprecise statistical
information, or as a revival of the seventeenth-century notion of
probability of testimony, as distinct from chance, due to Jakob Bernoulli
and later Lambert. The survey places Shafer's book in the second tradition.
& p.~311 & historical classification, secondhand, through
\cite{R47Shafer1978}; it concerns \cite{R47Dempster1967} and
\cite{R47Shafer1976} &
That Shafer's mathematics coincides with Dempster's is not shown here. The
classification is about lineage, not about which interpretation is correct. \\
The main problem for the seventeenth- and eighteenth-century pioneers was
merging testimonies. They proposed special cases of the rule now named
after Dempster, and they already suggested the reinforcement value
\(p_1+p_2-p_1p_2\). & p.~312 & historical attribution & The pioneers' own
rules are not reconstructed, and no claim is made that they assumed
independent sources. The modern rule and its independence hypothesis are
\cref{def:p15-r47-combination,thm:p15-r47-combination}. \\
One-sided betting, in which a buying price need not also be a selling
price, was pioneered by Smith (1961), Williams (1975) and Giles (1982), and
Walley (1991) developed it more extensively. & p.~312 & historical attribution
about \cite{R47Smith1961}, \cite{R47Williams2007}, \cite{R47Giles1982} and
\cite{R47Walley1991} & Williams's 1975 report is cited through its 2007
publication. Priority among these authors is not examined here. \\
Ellsberg's urn gambles: when people choose between gambles on a ball drawn
from an urn whose composition is poorly known, many experiments show
systematic violations of Savage's axioms, especially the sure-thing
principle, which the survey attributes to pessimism about rewarding
events of unknown probability. & p.~311 & report of an empirical
finding, with an explanatory hypothesis & The survey gives no task protocol,
sample, population or uncertainty estimate. The original article's urns and
the status of its observations are in
\cref{def:p17-ellsberg-three-colors,thm:p17-ellsberg-three-colors}
and the discussion following them; the decision-rule analysis is
\cref{thm:p15-w02b-p15-032-ellsberg}. \\
Kyburg used the lottery paradox to argue that accepted beliefs need not
be deductively closed, and with Teng built a logic for that view, the
logic of risky knowledge. & pp.~308, 310 & historical attribution about
\cite{R47Kyburg1961}, \cite{R47Kyburg1997} and \cite{R47KyburgTeng2012} &
That logic is not presented here. The lottery calculation is in
\cref{thm:p15-r47-big-stepped}, and the alternative of restricting the
probabilities is \cref{thm:p15-w02b-p15-032-big-stepped}. \\
\bottomrule\end{longtable}

\section{The uniqueness principle and the fair-price route}
\label{sec:p15-w02b-p15-032-uniqueness}

The survey states the Bayesian uniqueness principle as a starting point
for every kind of uncertainty modelling: any state of knowledge can be
represented by one probability distribution. It cites Lindley as one
statement of the principle \cite[p.~310]{DuboisPrade2021RepresentingBelief}
\cite{R47Lindley1982}. This is a normative and methodological thesis. Its
formal support in the survey is de Finetti's procedure. If an agent assigns
buying prices to all outcomes of a game of chance, and the same prices
are also selling prices, coherence forces one distribution. The survey
then argues that the resulting single measure cannot record the
difference between incomplete information and genuine randomness, citing
Shafer's 1976 objections \cite[p.~310]{DuboisPrade2021RepresentingBelief}.

\begin{hbdefinition}{def:p15-w02b-p15-032-two-sided}{Two-sided and one-sided event prices}
Let \(q:2^S\to\mathbb R\) give a price for every unit event ticket
\(\mathbf 1_A\). The schedule is \emph{two-sided} when the agent will buy
\(\mathbf 1_A\) at \(q(A)\) and also sell it at \(q(A)\); this is the
reversible contract model of \cref{def:p17-finite-price-coherence}.
The schedule is \emph{one-sided} when \(q(A)\) is only an acceptable
buying price. Write \(\Delta(S)\) for the probability vectors on \(S\).
Each price imposes a constraint on a probability \(P\) taken as a fair
valuation: buying at \(q(A)\) requires \(P(A)\geq q(A)\), and selling at
\(q(A)\) requires \(P(A)\leq q(A)\). The \emph{compatible set} of the
schedule is the set of \(P\in\Delta(S)\) meeting every constraint the
schedule imposes.
\end{hbdefinition}

\begin{hbtheorem}{thm:p15-w02b-p15-032-unique}{proved=annex:p15-w02b-p15-032-unique}{Fair prices force one distribution; one-sided prices need not}
\begin{enumerate}
\item If a two-sided schedule on all events avoids strict sure loss, then
  \(q\) is a probability and its compatible set is exactly \(\{q\}\).
\item If, in addition, the singleton prices on \(S\) are all equal, then
  \(q\) is the uniform distribution.
\item The one-sided vacuous schedule, \(q(S)=1\) and \(q(A)=0\) for
  \(A\neq S\), avoids sure loss. Its compatible set is all of
  \(\Delta(S)\), and every price in it is the lower envelope of that set.
\end{enumerate}
\end{hbtheorem}

Part~1 is the formal content of the survey's claim that coherence forces
one distribution. It depends on the two-sided commitment, which
\cref{sec:p15-rat41-probabilism-arguments} lists among the contested
premises of the Dutch book argument. Part~2 is the survey's die. An agent
who knows the die is fair and an agent with no information about it, if
both must quote two-sided symmetric prices, quote the same uniform vector.
Part~3 shows that one-sided prices keep the second agent's state distinct.
The vacuous schedule is the lower envelope of every distribution, while
the fair-die schedule is a single point. Neither part shows that the
first agent is irrational or that the second is. The formal results only
locate the premise from which uniqueness follows.

\section{Cox's postulates and what they presuppose}
\label{sec:p15-w02b-p15-032-cox}

The survey's version of Cox's postulates and regularity condition is
\cref{def:p15-r47-cox}. Cox claimed that a conditional confidence function
meeting them is isomorphic, by a rescaling, to conditional probability.
The survey reports three reservations about the proof
\cite[p.~310]{DuboisPrade2021RepresentingBelief}. Cox's original proof was
faulty. Paris gave another proof under a weaker regularity condition, in
which it suffices that \(F\) be strictly increasing in each argument, but he
also needed a density condition that requires an infinite setting
\cite{R47Paris1994}. Halpern showed that the result fails on finite
spaces and that Cox's original conditions do not suffice even in the
infinite setting \cite{R47Halpern1999a,R47Halpern1999b}. None of the three
works is read here. Their statements stay attributed to the survey, and
none of them is proved below.

The survey's further objections do not depend on those technical points.
Postulate~(i) looks natural only to someone who already takes the form of
Bayes conditioning for granted. Postulate~(ii) requires self-duality, which
excludes representations of partial ignorance. The survey concludes that
these points seriously weaken the claimed universality of Cox's result
\cite[p.~310]{DuboisPrade2021RepresentingBelief}. The following result
makes both objections precise. Here each \(g(\cdot\mid B)\) is assumed to
be a confidence measure in the sense of \cref{def:p15-r47-capacity}. That
assumption is the survey's own normalization, not one of Cox's postulates.

\begin{hbtheorem}{thm:p15-w02b-p15-032-cox}{proved=annex:p15-w02b-p15-032-cox}{The product form of postulate (i) and the ignorance exclusion of postulate (ii)}
\begin{enumerate}
\item For a probability \(P\) on \(S\), put \(g(A\mid B)=P(A\cap B)/P(B)\)
  when \(P(B)>0\). Then postulate~(i) holds with \(F(x,y)=xy\) whenever
  \(P(C\cap B)>0\), and postulate~(ii) holds with \(n(x)=1-x\).
\item Suppose each \(g(\cdot\mid B)\) satisfies \(g(\varnothing\mid B)=0\)
  and \(g(S\mid B)=1\), and postulate~(ii) holds with some function \(n\).
  Then \(n(0)=1\) and \(n(1)=0\). Consequently \(g(A\mid B)=0\) implies
  \(g(A^c\mid B)=1\), and no event and its complement both receive
  confidence zero in the same context.
\item Let \(N\) be a normalized necessity measure on finite \(S\) with an
  event \(A\) such that \(N(A)=N(A^c)=0\). Then \(g(\cdot\mid S)=N\) satisfies
  no postulate of form~(ii) at \(B=S\). For example, on \(S=\{1,2,3\}\)
  with \(\pi=(1,1,0)\), \(N(\{1\})=N(\{3\})=0\) while
  \(N(\{2,3\})=0\) and \(N(\{1,2\})=1\).
\end{enumerate}
\end{hbtheorem}

Part~1 shows that postulate~(i) generalizes the product rule
\(P(A\cap C\mid B)=P(A\mid C\cap B)P(C\mid B)\). That is the sense in which
it builds in the form of Bayes conditioning. The survey's point is a
normative one: whoever does not already accept that form has been given
no independent reason to accept~(i). Parts~2 and~3 make the self-duality
objection exact. In possibility theory the pair \(N(A)=N(A^c)=0\) is the
representation of complete ignorance about \(A\), and
\cref{thm:p15-r47-threshold} shows that \(\min\{N(A),N(A^c)\}=0\) always holds.
Postulate~(ii) excludes precisely that state. So Cox's argument, even with
a correct proof, would compare probability only with confidence functions
that already exclude this kind of ignorance.

\section{Superadditivity: two disjoint events against every order}
\label{sec:p15-w02b-p15-032-superadditivity}

The survey separates belief functions from general coherent lower
probabilities by an order of superadditivity
\cite[p.~313]{DuboisPrade2021RepresentingBelief}. Belief functions satisfy
the inclusion--exclusion inequality for every finite family of events; the
proof is \cref{thm:p15-r47-focal}. A lower probability obtained from an
arbitrary credal set is guaranteed only the two-event inequality for
disjoint events.

\begin{hblemma}{thm:p15-w02b-p15-032-disjoint}{proved=annex:p15-w02b-p15-032-disjoint}{Every lower envelope is superadditive on disjoint pairs}
Let \(\mathcal P\subseteq\Delta(S)\) be nonempty and put
\(P_*(A)=\inf_{P\in\mathcal P}P(A)\) and \(P^*(A)=1-P_*(A^c)\). Then
\(P_*(\varnothing)=0\), \(P_*(S)=1\), \(P_*\) is monotone, and
\[
 P_*(A\cup B)\geq P_*(A)+P_*(B)\qquad(A\cap B=\varnothing).
\]
Moreover \(P_*(A)\leq P(A)\leq P^*(A)\) for every \(P\in\mathcal P\).
\end{hblemma}

The inequality cannot be strengthened to the order-two inequality
\(P_*(A\cup B)+P_*(A\cap B)\geq P_*(A)+P_*(B)\) for overlapping events. The
two-point credal set of \cref{thm:p14-evidence-measure-separation} violates
it, so that lower envelope is not a belief function. The survey adds that
the M\"obius transform of such a \(P_*\) exists and is unique, but need not be
nonnegative. It credits this to Chateauneuf and Jaffray
\cite{R47ChateauneufJaffray1989}. Existence, uniqueness and the
nonnegativity test are \cref{thm:p15-r47-mobius}. In that counterexample the
negative mass is forced: if all inverse masses were nonnegative, the envelope
would be a belief function and would satisfy the order-two inequality.

The survey also says that the interval \([P_*(A),P^*(A)]\) is not meant to
contain an unknown subjective or objective probability of \(A\). On that
reading, degrees of belief are themselves precise numbers: they are the
coherent lower probabilities \cite[p.~313]{DuboisPrade2021RepresentingBelief}.
The last inequality of the lemma is a fact about the members of
\(\mathcal P\). It does not say that one member is the true one.

\section{Two decision rules for lower previsions}
\label{sec:p15-w02b-p15-032-decision}

The survey derives a decision rule from desirability. A gamble \(f\) is
preferred to \(g\) when \(f-g\) is desirable, that is, when the lower
expectation of \(f-g\) over the credal set is positive. This gives a
partial order, and it implies \(LP(f)\geq LP(g)\). The survey then contrasts
two rules. Comparing lower previsions, \(LP(f)\geq LP(g)\), solves the
Ellsberg paradox. The rule \(LP(f-g)\geq0\) satisfies Savage's sure-thing
principle \cite[p.~313]{DuboisPrade2021RepresentingBelief}. The survey writes
the second rule with a weak inequality, although desirability was stated with
a positive lower expectation. Both versions are recorded below.

\begin{hbdefinition}{def:p15-w02b-p15-032-rules}{Maxmin comparison and difference comparison}
Gambles are real vectors \(f\in\mathbb R^S\), read as utilities in each
state. Fix a nonempty credal set \(\mathcal P\subseteq\Delta(S)\) and put
\(LP(f)=\inf_{P\in\mathcal P}E_P(f)\). Define
\[
 f\succeq_M g\iff LP(f)\geq LP(g),\qquad
 f\succeq_D g\iff LP(f-g)\geq0,\qquad
 f\succ_D^{+} g\iff LP(f-g)>0.
\]
The first is \emph{maxmin comparison}, the second \emph{difference
comparison}, and the third \emph{strict desirability} of the difference.
For \(E\subseteq S\) and gambles \(f,h\), the pasted gamble \(f_Eh\) equals
\(f\) on \(E\) and \(h\) on \(S\setminus E\). A relation \(\succeq\)
satisfies the \emph{pasting condition} when
\(f_Eh\succeq g_Eh\iff f_Eh'\succeq g_Eh'\) for all \(f,g,h,h'\) and
\(E\). This is the sure-thing condition of
\cref{def:p17-ambiguity-acts}, stated here without requiring the relation
to be complete.
\end{hbdefinition}

\begin{hbtheorem}{thm:p15-w02b-p15-032-ellsberg}{proved=annex:p15-w02b-p15-032-ellsberg}{The two rules, the sure-thing condition and Ellsberg's three-colour urn}
\begin{enumerate}
\item If \(f\succeq_D g\) then \(f\succeq_M g\), and if \(f\succ_D^{+}g\)
  then \(LP(f)>LP(g)\). The relation \(\succeq_D\) is reflexive and
  transitive.
\item \(f\succeq_D g\) holds exactly when \(E_P(f)\geq E_P(g)\) for every
  \(P\in\mathcal P\). This is the robust dominance relation of
  \cref{def:p17-robust-dominance}, and no compactness of \(\mathcal P\)
  is needed.
\item \(\succeq_D\) and \(\succ_D^{+}\) satisfy the pasting condition.
\item Take the three-colour urn of \cref{def:p17-ellsberg-three-colors}
  with \(u(H)=1\), \(u(L)=0\), states in the order \((R,B,Y)\), and
  \[
   \mathcal P_E=\{(1/3,b,2/3-b):0\leq b\leq2/3\}.
  \]
  Then \(LP(F_1,F_2,F_3,F_4)=(1/3,0,1/3,2/3)\). So \(\succeq_M\) gives
  \(F_1\succ_M F_2\) and \(F_4\succ_M F_3\), which is Ellsberg's principal
  pattern, and it fails the pasting condition. Under \(\succeq_D\) both
  pairs are incomparable. Since \(F_1-F_2=F_3-F_4\), the rule \(\succeq_D\)
  treats the two pairs alike, as the pasting condition requires.
\item The basic assignment \(m(\{R\})=1/3\), \(m(\{B,Y\})=2/3\) has
  \(\mathcal P_E\) as the set of probabilities dominating its belief
  function. Its pignistic probability is uniform, so expected utility
  under it gives \(F_1\sim F_2\) and \(F_3\sim F_4\), not Ellsberg's pattern.
  The vacuous assignment \(m(S)=1\) also has the uniform pignistic
  probability.
\end{enumerate}
\end{hbtheorem}

\begin{hbexample}{def:p15-w02b-p15-032-urn-example}{Reading the urn calculation}
In part~4 the known red third fixes \(LP(F_1)=1/3\) and \(LP(F_4)=2/3\).
Neither bet depends on how the sixty other balls divide. The bets \(F_2\)
and \(F_3\) do depend on that division, and the worst case is taken. The
same values are \cref{thm:p17-maxmin-urns}, there computed from two extreme
compositions. Part~5 shows why the survey distinguishes Smets and Kennes's
pignistic rule from lower-envelope reasoning. A belief function that
represents the urn exactly, once projected to one probability, loses the
asymmetry that the maxmin rule uses. The uniform value for the vacuous
assignment is the sense in which the survey calls the pignistic
probability a generalization of Laplace's principle of insufficient reason
\cite[p.~312]{DuboisPrade2021RepresentingBelief}.
\end{hbexample}

``Solves'' in the survey means that maxmin comparison can represent the
reported pattern; that is the formal content of part~4. It is not
evidence that the pattern is rational. Nor is it evidence about how often
the pattern occurs, which the survey does not measure. The choice between
the rules is a normative trade-off. Maxmin comparison is complete and
represents the pattern, but fails the pasting condition. Difference
comparison keeps the pasting condition and transitivity but leaves some
pairs incomparable. Part~5 involves no measured choices. It is a
computation for one stipulated representation of the urn.

\section{Acceptance, the minimal epistemic logic and the bounded chain}
\label{sec:p15-w02b-p15-032-mel}

\Cref{def:p15-r47-mel} gives the language and models of the minimal
epistemic logic MEL. The survey says what its two modal axioms contribute
\cite[p.~309]{DuboisPrade2021RepresentingBelief}. Axiom~K makes
\(\Box(p\land q)\) equivalent to \(\Box p\land\Box q\), which is the
characteristic axiom of necessity measures. Axiom~D says that \(N(A)=1\)
implies \(\Pi(A)=1\); it secures consistency, since models of MEL are
nonempty subsets of \(S\). In GPL axiom~D takes the stronger form
\(\Box^\lambda p\Rightarrow\Diamond^1p\). The survey also records the KD45
axioms with an accessibility-relation semantics, in which \(\Box p\) holds at
\(s\) when every state accessible from \(s\) satisfies \(p\). It argues that
the introspection axioms, which it labels S4 and S5, mainly serve to reduce
nested modalities.
Evaluating belief at an epistemic state is more natural than evaluating it
at a real state of affairs. Those KD45 statements, and the
handbook's multi-agent KD45 system to which they correspond, are discussed
in \cref{sec:p15-rationality-4-7}.

\begin{hbtheorem}{thm:p15-w02b-p15-032-mel}{proved=annex:p15-w02b-p15-032-mel}{The set-function content of axioms K and D}
For \(E\subseteq S\), define \(N_E(A)=1\) if \(E\subseteq A\) and
\(N_E(A)=0\) otherwise, and \(\Pi_E(A)=1-N_E(A^c)\). Read \(\Box p\) at \(E\) as
\(N_E([p])=1\).
\begin{enumerate}
\item For every \(E\), including \(E=\varnothing\), the K-instance
  \(\Box(p\land q)\leftrightarrow(\Box p\land\Box q)\) and the schema
  \(\Box(p\Rightarrow q)\Rightarrow(\Box p\Rightarrow\Box q)\) hold. Thus
  \(N_E(A\cap B)=\min\{N_E(A),N_E(B)\}\).
\item The schema \(\Box p\Rightarrow\neg\Box\neg p\) holds at \(E\) for every
  \(p\) exactly when \(E\neq\varnothing\). Equivalently, \(N_E(A)=1\) implies
  \(\Pi_E(A)=1\) for every \(A\), and \(N_E(\varnothing)=0\).
\item Hintikka's ignorance about \(p\),
  \(\neg\Box p\land\neg\Box\neg p\), holds at a nonempty \(E\) exactly when
  \(E\) meets both \([p]\) and \([\neg p]\).
\item For a normalized possibility distribution with necessity \(N\) and
  possibility \(\Pi\), \(N(A)>0\) implies \(\Pi(A)=1\). Hence the graded
  form \(\Box^\lambda p\Rightarrow\Diamond^1p\) holds for every positive grade.
\end{enumerate}
\end{hbtheorem}

The survey goes on to claim that qualitative possibility and necessity
functions valued in a bounded chain give a unified framework for three
traditions. The first is modal logic, the second the set-function
tradition, and the third truth-functional many-valued logic, whose
representing power the survey describes as severely limited
\cite{R47CiucciDubois2013}. On this account accepted beliefs are
defeasible. The survey claims close links with nonmonotonic reasoning,
counterfactual reasoning and belief revision. It cites Dubois, Fargier and
Prade (2004) and refers to the handbook chapters by Rott and by Starr
\cite[p.~308]{DuboisPrade2021RepresentingBelief}
\cite{R47DuboisFargierPrade2004}. It ends its account of the graded logic
by calling GPL the logic of qualitative Baconian probabilities
\cite[p.~310]{DuboisPrade2021RepresentingBelief}. These are interpretive
claims about a research programme. What can be checked from the
definitions is stated next.

\begin{hbdefinition}{def:p15-w02b-p15-032-chain}{Chain-valued necessity and conditional acceptance}
A \emph{bounded chain} is a totally ordered set \(L\) with least element
\(0_L\) and greatest element \(1_L\). A \emph{chain-valued necessity} on
\(S\) is a map \(N:2^S\to L\) with \(N(\varnothing)=0_L\),
\(N(S)=1_L\) and \(N(A\cap B)=\min\{N(A),N(B)\}\). For a possibility
distribution \(\pi\) with values in \([0,1]\) and maximum one, say that
\(q\) is \emph{accepted in context} \(p\) when
\(\Pi([p\land q])>\Pi([p\land\neg q])\).
\end{hbdefinition}

\begin{hbtheorem}{thm:p15-w02b-p15-032-chain}{proved=annex:p15-w02b-p15-032-chain}{Chain representation, the conditional encoding and defeasible acceptance}
\begin{enumerate}
\item On finite \(S\), a chain-valued necessity satisfies
  \(N(A)=\min_{s\notin A}N(S\setminus\{s\})\), with the empty minimum
  \(1_L\), and some \(s\) has \(N(S\setminus\{s\})=0_L\). For
  \(L=\{0,1\}\) these are exactly the maps \(N_E\) of
  \cref{thm:p15-w02b-p15-032-mel} with \(E\neq\varnothing\).
\item \(N([p\Rightarrow q])>N([p\Rightarrow\neg q])\) holds exactly when \(q\)
  is accepted in context \(p\). This is the constraint by which the survey
  encodes the System~P conditional \cite[p.~310]{DuboisPrade2021RepresentingBelief}.
\item For each context \(p\) with \(\Pi([p])>0\), the formulas accepted in
  context \(p\) are closed under conjunction and do not include both a
  formula and its negation.
\item Acceptance in context is not monotone in the context. On the eight
  valuations of atoms \(b,g,f\), let \(\pi\) take the value \(1\) on
  \(b\bar gf\) and \(\bar b\bar g\bar f\), \(1/2\) on \(bg\bar f\), \(1/4\)
  on \(bgf\), \(b\bar g\bar f\) and \(\bar b\bar gf\), and \(0\) on
  \(\bar bg\bar f\) and \(\bar bgf\). Then \(f\) is accepted in context
  \(b\), and \(\neg f\) is accepted in context \(b\land g\).
\end{enumerate}
\end{hbtheorem}

Part~1 covers the survey's three cases with one representation: the
Boolean chain of MEL, the finite grades of GPL in \cref{def:p15-r47-gpl},
and the unit interval. Part~3 is the contextual analogue of adjunction.
Part~4 is a constructed illustration, not an example from the survey. If
\(b\) is read as ``bird'', \(g\) as ``penguin'' and \(f\) as ``flies'', the
example shows a defeasible conclusion. The conclusion is withdrawn when
the context is strengthened, while consequence between sets of models
remains monotone, as \cref{sec:p15-rationality-4-7} explains.
The grades used are \(\{0,1/4,1/2,3/4,1\}\), a uniform grid of the kind in
\cref{def:p15-r47-gpl}. These parts do not establish the survey's links with
counterfactual reasoning or belief revision. In the handbook, revision is
the AGM theory of \cref{def:p08-agm}, and counterfactual conditionals are
the logics of \cref{tab:p08-conditional-logics}. The survey's claim that
qualitative possibility connects with both rests on the cited 2004 paper,
which is not read here. Ranking functions, which give another ordinal
model of defeasible acceptance with their own conditioning, are
\cref{sec:p15-rationality-5-3}.

\section{Which convex sets lower previsions represent}
\label{sec:p15-w02b-p15-032-convex}

The survey states that any convex set of probabilities can be represented
by lower previsions on some family of gambles
\cite[p.~312]{DuboisPrade2021RepresentingBelief}. For a closed set this is
true even on the family of all gambles. Without closure it fails.

\begin{hbtheorem}{thm:p15-w02b-p15-032-convex}{proved=annex:p15-w02b-p15-032-convex}{Closed convex credal sets are exactly recovered; others are not}
Let \(S\) be finite and \(K\subseteq\Delta(S)\) nonempty and convex, with
\(LP_K(f)=\inf_{P\in K}E_P(f)\) for \(f\in\mathbb R^S\). Put
\[
 \mathcal M(LP_K)=\{P\in\Delta(S):E_P(f)\geq LP_K(f)\text{ for all }f\}.
\]
Then \(\mathcal M(LP_K)\) is the topological closure of \(K\). Hence
\(\mathcal M(LP_K)=K\) exactly when \(K\) is closed. For \(S=\{1,2\}\)
and \(K=\{(t,1-t):0<t\leq1\}\), the set \(\mathcal M(LP_K)\) also contains
\((0,1)\).
\end{hbtheorem}

The general bounded-gamble form, with finitely additive probabilities and
the weak-star topology, is \cref{thm:p15-ip-representation}. There the
compatible set of a coherent functional on all gambles is its unique
closed convex representing set. The finite proof uses only a nearest-point
separation. The survey's ``some family of gambles'' therefore has to be
read with closure added, and on a restricted family of gambles even a
closed set need not be recovered, as \cref{sec:p15-r47-lower-previsions}
explains.

\section{Probabilities that preserve adjunction}
\label{sec:p15-w02b-p15-032-big-stepped}

The survey's adjunction rule at threshold one half requires, for every
context \(C\) and all events \(A,B\), that \(P(A\mid C)>P(A^c\mid C)\) and
\(P(B\mid C)>P(B^c\mid C)\) imply \(P(A\cap B\mid C)>P(A^c\cup B^c\mid C)\). It
says that the probabilities obeying this rule have masses
\(p_1>p_2>\cdots>p_n\) with \(p_i>p_{i+1}+\cdots+p_n\) for all \(i<n\)
\cite[p.~308]{DuboisPrade2021RepresentingBelief}. It credits these to
Benferhat, Dubois and Prade and to Snow
\cite{R47BenferhatDuboisPrade1999,R47Snow1999}. \Cref{thm:p15-r47-big-stepped}
proves that such masses are sufficient. The exact class is slightly larger.

\begin{hbtheorem}{thm:p15-w02b-p15-032-big-stepped}{proved=annex:p15-w02b-p15-032-big-stepped}{The exact finite class preserving adjunction at one half}
Let \(P\) be a probability on finite \(S\). Index its positive-mass states
so that \(p_1\geq p_2\geq\cdots\geq p_k>0\). Read the rule only for contexts
with \(P(C)>0\). Then \(P\) obeys the adjunction rule exactly when
\[
 p_i>\sum_{j=i+1}^{k}p_j\qquad\text{for every }i\leq k-2.
\]
States of mass zero are unrestricted. In particular the uniform law on two
states obeys the rule although \(p_1>p_2\) fails, and the uniform law on
three states violates it.
\end{hbtheorem}

The printed strict chain therefore describes the case without ties and
without null states. Equality is possible only between the two smallest
positive masses. The survey also says that the same kind of result holds
for any acceptance threshold. That further claim is not proved here.
It remains a claim attributed to the cited papers, which are not read
here. The interpretive conclusion, that accepted belief fits probability
best for distributions strongly biased toward one outcome, is the
survey's reading of the theorem.

\section{Relationships with the subjective probability and decision chapters}
\label{sec:p15-w02b-p15-032-relations}

The survey sends readers to H\'ajek and Staffel's chapter for subjective
probability \cite[p.~306]{DuboisPrade2021RepresentingBelief}; that chapter
is \cref{sec:p15-rationality-4-1}. The two treatments agree on the formal
core. Both rest probabilism's betting argument on the finite Dutch book
theorem \cref{thm:p17-finite-price-coherence}, and both note that
uniformity depends on how the possibilities are represented. Compare
\cref{thm:p15-rat41-indifference} with \cref{thm:p15-r47-ignorance}. They
differ in three material ways.
\begin{itemize}
\item \emph{Two senses of uniqueness.} H\'ajek and Staffel's uniqueness
  thesis, in \cref{sec:p15-rat41-nonideal}, says that a body of evidence
  permits only one rational credence function. That is a claim about
  rational permission. Lindley's principle, as Dubois and Prade report it,
  says that every state of knowledge can be represented by one
  distribution. That is a claim about representation. A permissivist can
  accept Lindley's principle, and a uniqueness theorist can reject it.
\item \emph{Response to imprecision.} H\'ajek and Staffel discuss interval
  credences and object that exact endpoints add artificial precision.
  Dubois and Prade's lower-prevision account answers in a different way.
  The endpoints are precise lower and upper prices, and they are not meant
  to bracket a hidden true probability
  (\cref{sec:p15-w02b-p15-032-superadditivity}). It is a reinterpretation
  of the objection, not a demonstration that the objection fails.
\item \emph{Where additivity is questioned.} H\'ajek and Staffel weigh
  three arguments for probabilism, including accuracy dominance, which
  Dubois and Prade do not discuss. Dubois and Prade instead question the
  two-sided price premise, the hypotheses of Cox's theorem and the
  adequacy of additivity for ignorance
  (\cref{thm:p15-w02b-p15-032-unique,thm:p15-w02b-p15-032-cox}). Neither
  treatment refutes the arguments considered by the other.
\end{itemize}

For qualitative decision rules the survey refers to Hill's chapter
\cite[p.~308]{DuboisPrade2021RepresentingBelief}; that chapter is
\cref{sec:p17-rationality-8-4}. Maxmin comparison is Hill's maximin
expected utility \(V_M\), and difference comparison is Hill's unanimity
relation, both in \cref{def:p17-rat84-prior-rules}. Hill's appraisal agrees
with \cref{thm:p15-w02b-p15-032-ellsberg}: maxmin keeps a weak order but
gives up the sure-thing principle, while unanimity keeps that principle
and transitivity but not completeness (\cref{sec:p17-rat84-choice-axioms}).
The material differences are these.
\begin{itemize}
\item \emph{Numerical against ordinal rules.} The rules Dubois and Prade
  name at this point, Wald's pessimistic and optimistic criteria extended
  to possibility distributions and their leximin and leximax refinements,
  are qualitative. They are credited to Dubois, Prade and Sabbadin (2001)
  and to Fargier and Sabbadin (2005)
  \cite{R47DuboisPradeSabbadin2001,R47FargierSabbadin2005}. Hill's criteria
  all use numerical expected utilities over priors, capacities or
  second-order probabilities. The survey prints no formula for the
  qualitative criteria, so they are not identified here with any of
  Hill's rules.
\item \emph{Capacities.} Hill's Choquet evaluation allows any normalized
  monotone capacity. A capacity with singleton values \(.75\) and \(.75\)
  on two states, as in \cref{thm:p17-rat84-choquet}, is not even
  superadditive on disjoint events, so it is not a lower probability by
  \cref{thm:p15-w02b-p15-032-disjoint}. Belief functions are convex
  capacities, being superadditive at every order. For them the Choquet
  value equals maxmin over the core by \cref{thm:p17-rat84-choquet}, and the
  core is the credal set used in part~5 of
  \cref{thm:p15-w02b-p15-032-ellsberg}.
\item \emph{Hurwicz-type rules.} Dubois and Prade mention Jaffray's
  axiomatized extension of Hurwicz's criterion to belief functions
  \cite{R47Jaffray1989}, and Hill presents \(\alpha\)-maxmin over a prior
  set. The survey does not state Jaffray's criterion, so the two are not
  claimed to coincide.
\end{itemize}

\section{Dependent sources for each claim}
\label{sec:p15-w02b-p15-032-sources}

\Cref{tab:p15-w02b-p15-032-claims} lists each proof or evidence claim that
the survey makes. For each one it gives how the claim stands in this
handbook: proved in this chapter, proved in another named result, or
attributed to a cited work that has not been read and whose proof is not
given here. Works in the last column are given by bibliographic identity
only; no page in them has been checked. Page numbers refer to the survey
\cite{DuboisPrade2021RepresentingBelief}.

\begin{longtable}{@{}p{.36\linewidth}p{.07\linewidth}p{.27\linewidth}p{.22\linewidth}@{}}
\caption{Claims of the 2021 survey and the result or source on which each depends.}
\label{tab:p15-w02b-p15-032-claims}\\
\toprule Claim & Page & Standing in this handbook & Dependent work \\
\midrule\endfirsthead
\caption[]{Claims of the 2021 survey and their dependencies (continued)}\\
\toprule Claim & Page & Standing in this handbook & Dependent work \\
\midrule\endhead
\midrule\multicolumn{4}{r}{\emph{continued on next page}}\\
\endfoot
\endlastfoot
Fair two-sided prices that avoid sure loss are additive & 306 &
proved: \cref{thm:p17-finite-price-coherence} & \cite{R47DeFinetti1937} \\
Coherence of fair prices on all outcomes forces one distribution & 310 &
proved: \cref{thm:p15-w02b-p15-032-unique} & \cite{R47Lindley1982} \\
Finite Boolean necessity has a nonempty epistemic state; every-threshold
adjunction is minitivity; \(N=\min\iota\), \(\Pi=\max\pi\) & 307 &
proved: \cref{thm:p15-r47-threshold} & \cite{R47DuboisPrade1988} \\
Ranking functions give a degree of potential surprise by a power of two &
307 & proved with corrected orientation: \cref{thm:p15-r47-ranking} &
\cite{R47Spohn1988} \\
Subjective probability is justified from first principles for decisions
under uncertainty & 307 & cited only for Savage's own derivation; the
postulates are stated in \cref{def:p17-sa-postulates}, the converse is
proved in \cref{thm:p17-sa-general-converse}, and the existence direction
is cited in \cref{thm:p17-sa-general-source} & \cite{R47Savage1954} \\
Possibility theory has a Savage-style justification on finite \(S\) &
308 & cited only & \cite{R47DuboisPradeSabbadin2001} \\
Leximin and leximax refinements are special forms of expected utility &
308 & cited only & \cite{R47FargierSabbadin2005} \\
Thresholded probability fails conjunctive closure (lottery) & 308 &
proved: \cref{thm:p15-r47-big-stepped} & \cite{R47Kyburg1961} \\
Adjunction at one half holds exactly for big-stepped probabilities & 308 &
proved in corrected form: \cref{thm:p15-w02b-p15-032-big-stepped} &
\cite{R47BenferhatDuboisPrade1999} \\
The same kind of result holds at any threshold & 308 & cited only &
\cite{R47BenferhatDuboisPrade1999}, \cite{R47Snow1999} \\
Bounded-chain possibility unifies the traditions, with defeasible accepted
beliefs & 308 & partly proved: \cref{thm:p15-w02b-p15-032-chain};
links to counterfactuals and revision cited only &
\cite{R47DuboisFargierPrade2004} \\
Kleene-style truth-functional logic severely limits representation & 308 &
cited only & \cite{R47CiucciDubois2013} \\
MEL: K gives minitivity, D gives \(N(A)=1\Rightarrow\Pi(A)=1\), models are
nonempty & 309 & proved: \cref{thm:p15-w02b-p15-032-mel} &
\cite{R47BanerjeeDubois2014} \\
GPL is sound and complete & 309 & cited only & \cite{R47DuboisPradeSchockaert2017} \\
GPL encodes relative belief \(N(p)>N(q)\) & 309 & proved:
\cref{thm:p15-r47-finite-modal} & \cite{R47DuboisPradeSchockaert2017} \\
GPL encodes System~P by \(N(p\Rightarrow q)>N(p\Rightarrow\neg q)\) & 309--310 &
the semantic equivalence is proved: \cref{thm:p15-w02b-p15-032-chain}; the
encoding theorem is cited only & \cite{R47DuboisPradeSchockaert2017},
\cite{KrausLehmannMagidor1990} \\
GPL encodes standard possibilistic logic and answer-set programs; standard
possibilistic logic is nonmonotonic & 309--310 & cited only &
\cite{R47DuboisLangPrade1994}, \cite{R47BenferhatDuboisPrade1998} \\
Cox's postulates yield conditional probability & 310 & cited only &
\cite{R47Cox1946} \\
Paris's proof under strict monotonicity of \(F\) and a density condition &
310 & cited only & \cite{R47Paris1994} \\
Cox's result fails on finite spaces; original conditions insufficient &
310 & cited only & \cite{R47Halpern1999a}, \cite{R47Halpern1999b} \\
Postulate~(i) presupposes Bayes conditioning; (ii) excludes partial
ignorance & 310 & proved in the precise form of
\cref{thm:p15-w02b-p15-032-cox} & none \\
Refinement and log-scale uniformity conflict & 311 & proved:
\cref{thm:p15-r47-ignorance} & \cite{R47Shafer1976} \\
People systematically violate Savage's axioms in urn gambles & 311 &
empirical report; no data in the survey; status of the original
observations in \cref{sec:p17-ambiguity-and-robust-choice} and
\cref{def:p17-ellsberg-three-colors} & \cite{Ellsberg1961Ambiguity} \\
Belief functions include probabilities and Boolean and graded necessities &
312 & proved: \cref{thm:p15-r47-focal} & \cite{R47Shafer1976} \\
\(\operatorname{Bel}(A)+\operatorname{Bel}(A^c)\leq1\) and
\(\operatorname{Pl}(A)=1-\operatorname{Bel}(A^c)\) & 312 & proved:
\cref{thm:p14-evidence-finite-measures} & \cite{R47Shafer1976} \\
Dempster's rule and the reinforcement value \(p_1+p_2-p_1p_2\) & 312 &
proved: \cref{thm:p15-r47-combination} & \cite{R47Dempster1967} \\
Only separable belief functions combine from simple supports; a
generalized decomposition & 312 & cited only & \cite{R47Shafer1976},
\cite{R47Smets1995} \\
Pignistic probability equals the Shapley value and generalizes Laplace's
principle & 312 & proved: \cref{thm:p15-r47-pignistic},
\cref{thm:p15-w02b-p15-032-ellsberg} & \cite{R47SmetsKennes1994},
\cite{R47Shapley1953} \\
Jaffray axiomatizes a Hurwicz-type criterion & 312 & cited only &
\cite{R47Jaffray1989} \\
Lower previsions are lower expectations over a credal set; \(UP(f)=-LP(-f)\);
empty credal set means sure loss; coherence means \(LP=bp\) & 312--313 &
proved: \cref{thm:p15-ip-representation,thm:p15-ip-dual} &
\cite{R47Walley1991} \\
Avoiding sure loss means \(UP(f)\geq LP(f)\) for all gambles \(f\) &
313 & the necessary direction is proved: \cref{thm:p15-ip-dual}; read as
sufficient for separately quoted buying and selling prices it fails:
\cref{thm:p15-r47-price-boundaries} & \cite{R47Walley1991} \\
Any convex set is represented by lower previsions & 312 & proved for
closed sets, refuted otherwise: \cref{thm:p15-w02b-p15-032-convex} &
\cite{R47Walley1991} \\
Belief functions are superadditive at every order, lower probabilities
only on disjoint pairs & 313 & proved: \cref{thm:p15-r47-focal},
\cref{thm:p15-w02b-p15-032-disjoint},
\cref{thm:p14-evidence-measure-separation} & none \\
The M\"obius transform of a lower probability exists, is unique, and may be
negative & 313 & proved: \cref{thm:p15-r47-mobius},
\cref{thm:p14-evidence-measure-separation} & \cite{R47ChateauneufJaffray1989} \\
Belief functions are the only rational degrees of belief under a
Walley-style betting framework & 313 & cited only &
\cite{R47KerkvlietMeester2018} \\
Difference comparison is a partial order implying \(LP(f)\geq LP(g)\);
maxmin solves Ellsberg; \(LP(f-g)\geq0\) satisfies the sure-thing principle &
313 & proved: \cref{thm:p15-w02b-p15-032-ellsberg} & none \\
Necessity measures are the belief functions obeying adjunction; the
Baconian condition matches overlapping focal sets & 313 & proved:
\cref{thm:p15-r47-focal} & none \\
Necessity measures are coherent lower probabilities with Giles's cautious
betting & 313 & proved: \cref{thm:p15-ip-possibility-bridge}; the betting
reading is cited only & \cite{R47Giles1982} \\
The least informative belief function inducing a given probability is a
necessity measure & 313 & cited only; a constructive inverse is proved in
\cref{thm:p15-r47-pignistic} & \cite{R47DuboisPradeSmets2008} \\
Ranking-induced necessities relate to infinitesimal probabilities & 313 &
cited only; see the qualification after \cref{thm:p15-r47-ranking} &
\cite{R47Spohn1990} \\
The only meaningful extension of likelihood to sets is the maximum & 314 &
proved under maxitivity: \cref{thm:p15-r47-likelihood} &
\cite{R47DuboisMoralPrade1997}, \cite{R47Dubois2006} \\
Argumentation as a route from conflicting evidence to the provable & 314 &
interpretive proposal; see \cref{sec:p18-rationality-5-5} & \cite{R47Haenni2009} \\
\bottomrule\end{longtable}

The table is the complete set of dependent works for this chapter. No
other source supports a claim made here.
\index{uniqueness principle!Lindley}\index{Cox's postulates}
\index{lower probability!superadditivity}\index{sure-thing principle!difference comparison}
\index{minimal epistemic logic}\index{big-stepped probability}
